\documentclass[11pt,a4paper]{article}
\usepackage[margin=18mm]{geometry}
\usepackage{amsmath,amssymb,booktabs,mathrsfs}
\usepackage[hidelinks]{hyperref}
\newcommand{\Ev}{E_\nu}
\newcommand{\Ep}{E'}
\newcommand{\cR}{\mathcal{R}}
\newcommand{\cH}{\mathcal{H}}
\newcommand{\Rb}[1]{R_{#1}}
\newcommand{\Hb}[1]{H_{#1}}
\newcommand{\Phit}{\Phi(2\,\mathrm{MeV})}
\newcommand{\Etwo}{2\,\mathrm{MeV}}
\setlength{\parskip}{3pt}
\title{\vspace{-12mm}What the code computes, in the notation of the paper\vspace{-2mm}}
\author{Ko Yasuda}
\date{August 19, 2026}
\begin{document}
\maketitle
\vspace{-9mm}
Equation numbers refer to arXiv:2608.16107v1. All quantities below are exactly those of the paper; only one
auxiliary function, $H_i(\Ev)$ in Eq.~(N1), is introduced because it is needed to state precisely what the code does.
Throughout, the split energy is $\Etwo$ (as in the paper and in the code).

\section*{1. Definitions taken from the paper}
\begin{itemize}\itemsep1pt
\item $\phi(\Ev)$: differential flux (Appendix~B); $\Phi(\Ev)=\int_{\Ev}^\infty d\tilde E_\nu\,\phi(\tilde E_\nu)$, Eq.~(2.3), so $\phi=-d\Phi/d\Ev$.
\item $\partial\cH/\partial\Ep(\Ev,\Ep)$: response function acting on $\phi$, Eqs.~(4.4)--(4.5), containing $\epsilon(\Ep)$, the resolution
$G(\Ep,E_R)$ and $\partial\sigma/\partial E_R$. $\partial\cR/\partial\Ep\equiv\partial_{\Ev}(\partial\cH/\partial\Ep)$, Eq.~(4.8), acting on $\Phi$.
\item $R_i(\Ev)\equiv R_{[E'_i,E'_{i+1}]}(\Ev)=\int_{E'_i}^{E'_{i+1}}d\Ep\,\partial\cR/\partial\Ep$, Eq.~(4.12), in the notation of Fig.~3;
$R_{ij}=\int_{\Ev^j}^{\Ev^{j+1}}d\Ev\,R_i(\Ev)$, Eq.~(5.3).
\item Ansatz: $\Phi(\Ev)=\Phi_j$ on the $j$th of the $N_{\rm int}$ intervals of $[\Ev^{\min},\Etwo]$, Eq.~(5.1).
\end{itemize}
The bin-integrated response acting on $\phi$ (the same integral as Eq.~(4.12) but of $\partial\cH/\partial\Ep$) is
\begin{equation}
\Hb{i}(\Ev)\equiv\int_{E'_i}^{E'_{i+1}}d\Ep\,\frac{\partial\cH}{\partial\Ep}(\Ev,\Ep),
\qquad
\Rb{i}(\Ev)=\frac{d\Hb{i}}{d\Ev},\qquad \Hb{i}(0)=0 .
\tag{N1}
\end{equation}
$H_i(0)=0$ for the reason given below Eq.~(4.7) ($E_R^{\max}\to0$). Hence $H_i(\Ev)=\int_0^{\Ev}dE\,R_i(E)$ and, with the
intervals of Eq.~(5.3),
\begin{equation}
\Hb{i}(\Etwo)=\int_{\Ev^{\min}}^{\Etwo}d\Ev\,\Rb{i}(\Ev)=\sum_{j=1}^{N_{\rm int}}R_{ij}.
\tag{N2}
\end{equation}
This answers the question about the boundary term: the function multiplying $\Phi(\Etwo)$ is $\partial\cH/\partial\Ep$ at
$\Ev=\Etwo$ integrated over the bin, i.e.\ $\Hb{i}(\Etwo)$, and by (N2) it is just the row sum of the response matrix.

\section*{2. What the code computes}
Rates per unit mass (kg$^{-1}$yr$^{-1}$); the exposure $\mathcal E$ multiplies everything as in Eqs.~(4.10), (5.4).
\begin{align}
N_i^{\rm tot}&=\int_0^\infty d\Ev\,\phi(\Ev)\,\Hb{i}(\Ev)
&&\text{Eqs.~(4.6),(4.10); black line of Fig.~2,}\tag{N3}\\
N_i^{<2}&=\int_0^{\Etwo}d\Ev\,\phi(\Ev)\,\Hb{i}(\Ev)
&&\text{Eq.~(4.6), upper limit $\Etwo$; cyan line of Fig.~2,}\tag{N4}\\
h_i^{\rm code}&=N_i^{\rm tot}-N_i^{<2}=\int_{\Etwo}^\infty d\Ev\,\phi(\Ev)\,\Hb{i}(\Ev)
&&\text{green line of Fig.~2 ($=h_i$ of Ref.~[6]).}\tag{N5}
\end{align}
$N_i^{\rm tot}$ agrees with Fig.~3 of Ref.~[6]; $h_i^{\rm code}$ is what the text calls the contribution of neutrinos with $\Ev>\Etwo$.
The mock data are $O_i=N_i^{\rm tot}+b_i$ (Poisson fluctuated per pseudo-experiment). The fit minimises exactly Eq.~(5.7),
\begin{equation}
\chi^2=\sum_{i=1}^{d}\frac{\big[\mathcal E\big(\sum_{j}R_{ij}\,x_j+h_i^{\rm code}+b_i\big)-\mathcal E O_i\big]^2}{\mathcal E O_i},
\qquad x_j\ge0,\ \ x_{j+1}\le x_j ,
\tag{N6}
\end{equation}
with the constraints of Eq.~(5.8), the same $R_{ij}$, the same $N_{\rm int}$, the same solver and the same Monte-Carlo band
construction as in Appendix~C. The only difference from Eq.~(5.4) is that $h_i$ is $h_i^{\rm code}$ of (N5), and the
fitted parameters are called $x_j$ here because, as shown next, they are not $\Phi_j$.

\section*{3. What the fitted parameters are}
Apply the integration by parts of Eq.~(4.7) to $N_i^{<2}$, i.e.\ with the upper limit $\Etwo$ instead of $\infty$:
\begin{equation}
N_i^{<2}=\Big[-\Phi(\Ev)\,\Hb{i}(\Ev)\Big]_0^{\Etwo}+\int_0^{\Etwo}d\Ev\,\Phi(\Ev)\,\Rb{i}(\Ev)
=-\Phit\,\Hb{i}(\Etwo)+\sum_j R_{ij}\Phi_j .
\tag{N7}
\end{equation}
The lower boundary term vanishes as in the paper. The upper one does not: $\Phit\neq0$ (Fig.~1). Using (N2),
\begin{equation}
\boxed{\;N_i^{<2}=\sum_j R_{ij}\,\big(\Phi_j-\Phit\big)\;}
\qquad\Longrightarrow\qquad
x_j=\Delta\Phi_j\equiv\Phi_j-\Phit .
\tag{N8}
\end{equation}
Equivalently, defining $h_i^{\rm paper}\equiv\int_{\Etwo}^\infty d\Ev\,\Phi(\Ev)\Rb{i}(\Ev)$ (the $\Phi$-weighted response above $\Etwo$),
\begin{equation}
h_i^{\rm code}=h_i^{\rm paper}+\Phit\sum_j R_{ij},
\qquad
N_i^{\rm tot}=\underbrace{\sum_j R_{ij}\Phi_j+h_i^{\rm paper}}_{\text{Eq.~(5.4)}}
=\underbrace{\sum_j R_{ij}\Delta\Phi_j+h_i^{\rm code}}_{\text{code}} .
\tag{N9}
\end{equation}
Both decompositions are exact and give the same total; they differ in whether the term $\Phit\sum_j R_{ij}$ is
counted as signal or as ``high-energy'' contribution. Eq.~(5.4) with $\Phi_j$ requires $h_i^{\rm paper}$; the code
(and the definition of $h_i$ in the text and in Fig.~2 as the rate from neutrinos above $\Etwo$) uses $h_i^{\rm code}$,
which forces the fitted parameters to be $\Delta\Phi_j$. I have checked (N7)--(N9) numerically with the response table
of the code (better than 1\%).

\section*{4. Why $\Delta\Phi$ appeared; is it better?}
It was not a choice and it is not better. I computed $h_i$ as the physical rate from neutrinos above $\Etwo$, following
Ref.~[6] and Fig.~2, and this automatically turns the fitted parameters into $\Delta\Phi_j$. There is no advantage: the
total rate, $\chi^2$ and constraints are the same, but what comes out is $\Phi(\Ev)-\Phit$, which is why the
best-fit and the theory curves in Fig.~7 lie below Fig.~1 by $\Phit\simeq4.6\times10^{11}\,$cm$^{-2}$s$^{-1}$
($24\%$ of $\Phi(0.18\,\mathrm{MeV})$; solid curve of Fig.~1). The paper is written for $\Phi$, and that is what we want.

\section*{5. Size of the term (1 eV threshold; kg$^{-1}$yr$^{-1}$ per bin; solid flux of Fig.~1)}
\begin{center}\small
\begin{tabular}{r rrrrrr}
\toprule
bin [eV] & $N_i^{<2}=\sum_jR_{ij}\Delta\Phi_j$ & $\Phit\sum_jR_{ij}$ & $\sum_jR_{ij}\Phi_j$ & $h_i^{\rm paper}$ & $h_i^{\rm code}$ & $N_i^{\rm tot}$\\
\midrule
   1--3 &   232.5 &    96.9 &   329.4 &     0.8 &    97.7 &   330.2 \\
 11--13 &   121.9 &    92.3 &   214.2 &     4.9 &    97.3 &   219.2 \\
 21--23 &    76.1 &    83.7 &   159.8 &     9.1 &    92.8 &   168.9 \\
 41--43 &    33.0 &    66.5 &    99.5 &    17.3 &    83.8 &   116.8 \\
 51--56 &    52.0 &   141.5 &   193.5 &    55.1 &   196.6 &   248.6 \\
81--120 &    32.7 &   316.7 &   349.3 &   806.5 &  1123.1 &  1155.8 \\
\bottomrule
\end{tabular}
\end{center}
The term $\Phit\sum_jR_{ij}$ is comparable to or larger than $N_i^{<2}$ in every bin, so the difference between fitting
$\Phi_j$ and $\Delta\Phi_j$ is large.

\section*{6. What has to change}
Two self-consistent options:
\begin{itemize}\itemsep1pt
\item[(a)] \textbf{Fit $\Phi$, as written in Eq.~(5.4).} Signal $=\sum_jR_{ij}\Phi_j=N_i^{<2}+\Phit\sum_jR_{ij}$,
background $=h_i^{\rm paper}$. All fits and bands must be redone (both thresholds); Fig.~2 (green and cyan lines) and the sentence
defining $h_i$ must be changed ($h_i$ is then the $\Phi$-weighted response above $\Etwo$, not the number of events from
neutrinos above $\Etwo$; at low $\Ep$ it is tiny, see the table). The comparison with Ref.~[6] must be redone.
\item[(b)] \textbf{Keep the numerics, rewrite in $\Delta\Phi$.} Eqs.~(5.1), (5.2), (5.4), (5.7), (5.8) and the vertical axes of the
flux figures in terms of $\Delta\Phi_j=\Phi_j-\Phit$ (constraints $\Delta\Phi_j\ge0$, non-increasing, hold as well); $h_i$, Fig.~2 and all
fits unchanged; the theory curves shifted by $\Phit$.
\end{itemize}
Since the paper is written for $\Phi$, I propose (a). I can start immediately; the code changes are only the replacement of the
signal and $h_i$ vectors, but the confidence bands have to be recomputed.

\vspace{2mm}\footnotesize
Note on inputs: the code implements the resolution as a box of half-width 1~eV (the paper states this is equivalent to the
Gaussian of Eq.~(3.3)) and $\epsilon(\Ep)=\theta(\Ep-E'_{\rm thr})$; the numbers in the table use the response table of the code and the
solid flux of Fig.~1.
\end{document}
